\scanpage{120}
\mspage{4}
% archive sheet number: 119
% margin note written slanted, rising to the right
\margin{\uline{Cor} (de $\mathrm{Cart}(3)$ et $\mathrm{Cart}(2)$) \unsure{tt} carré isomorphe à un carré $h$-cartésien est $h$-cartésien.}
des carrés $h$-cartésiens (\unsure{savoir} la face inférieure et les quatre faces latérales), il en est de même de la face supérieure.
\deleted{\uline{NB} En fait, il suffit de supposer que}

Ces axiomes sont dans la nature d'axiomes de stabilité, nous donnant des critères \unsure{suffisants} pour que certains carrés \unsure{sont|soient} $h$-cartésiens, si telles conditions sont vérifiées (p. ex. que certains autres carrés sont déjà $h$-cartésiens). J'ai l'impression d'avoir \deleted{\unsure{ici}} épuisé ici les propriétés \unsure{pertinentes} de ce type (i.e. que les autres devraient pouvoir s'en déduire). À présent, je m'occupe des propriétés d'existence et d'unicité des carrés $h$-cartésiens.

% vertical emphasis stroke beside Cart(4)
\margin{$\vert$}
\uline{Cart(\unsure{4|5})} \inserted{(Axiome des produits fibrés)} Tout diagramme
\diagram{\begin{tikzcd}[ampersand replacement=\&]
X_1 \arrow[d] \& \\
\unsure{X_0} \& X_2 \arrow[l]
\end{tikzcd}}
dans $\unsure{\mathscr H|\mathcal H|H}$ peut se \unsure{complèter|compléter} en un carré $h$-cartésien
\diagram{\begin{tikzcd}[ampersand replacement=\&]
X_1 \arrow[d] \& X_3 \arrow[l] \arrow[d] \\
X_0 \& X_2 \arrow[l]
\end{tikzcd}}
.

L'axiome suivant \inserted{(cf \unsure{son} corollaire)} \unsure{assurera|assure} que le carré $h$-cartésien en question est \inserted{\unsure{intrinsèquement}} \uline{unique à isom} près - mais on fera attention qu'il n'est en général pas unique ``à isom unique près''.

\uline{Cart(5)} \inserted{(Axiome d'exactitude)} Considérons deux \deleted{\unsure{diagr.}} carrés $h$-cartésiens
\diagram{\begin{tikzcd}[ampersand replacement=\&]
X=\Biggl( \& X_1 \arrow[d] \& X_3 \arrow[l] \arrow[d] \& \Biggr) \\
\& X_0 \& X_2 \arrow[l] \&
\end{tikzcd}}
et
\diagram{\begin{tikzcd}[ampersand replacement=\&]
Y=\Biggl( \& Y_1 \arrow[d] \& \unsure{X_3|Y_3} \arrow[l] \arrow[d] \& \Biggr) \\
\& Y_0 \& Y_2 \arrow[l] \&
\end{tikzcd}}
, et un hom. \inserted{$p=(p_0,p_1,p_2)$} entre les diagrammes de type
\diagram{\begin{tikzcd}[ampersand replacement=\&]
\unsure{\Phi|\mathscr Q}= \& 1 \arrow[d] \& \\
\& 0 \& 2 \arrow[l]
\end{tikzcd}}
associés (cf diagrammes traits pleins ci-dessous) :
